\chapter{Gout}

\section*{Definition and Overview}
Gout is a form of inflammatory arthritis caused by high levels of uric acid in the blood, which can form sharp crystals in the joints and trigger sudden, severe attacks of pain, redness, and swelling. It most commonly affects the big toe but can affect the ankles, knees, wrists, fingers, and elbows. Gout attacks typically come on suddenly, often at night, peak within hours, and resolve within a week or two even without treatment. Between attacks, uric acid can continue to deposit in joints, causing long-term damage, and repeated attacks can lead to chronic gouty arthritis and tophi (crystal deposits under the skin). Gout is highly treatable with medicines that relieve attacks and lower uric acid levels.

\section*{Epidemiology}
Gout is the most common form of inflammatory arthritis in Australia, affecting around 4--6\% of adults, with rates increasing with age. It is several times more common in men than women, and in women it most often begins after menopause. Prevalence has risen over recent decades, partly in line with increasing rates of obesity, high blood pressure, and chronic kidney disease, and the increasing use of medicines such as diuretics that raise uric acid. Gout is strongly associated with diets high in red meat, seafood, and alcohol. Without treatment, many people have recurrent attacks over many years.

\section*{Aetiology and Pathophysiology}
Gout is caused by an excess of uric acid that crystallises in the joints.

\begin{itemize}
    \item \textbf{Uric acid production:} uric acid is a breakdown product of purines, substances found in many foods and produced by the body
    \item \textbf{Under-excretion:} in most people with gout, the kidneys do not excrete uric acid efficiently enough
    \item \textbf{Hyperuricaemia:} when uric acid levels rise above the saturation point, monosodium urate crystals form in joints and soft tissues
    \item \textbf{Crystal-induced inflammation:} the crystals trigger an intense inflammatory response, causing the classic acute attack
    \item \textbf{Triggers:} alcohol, dehydration, a rich meal, injury, surgery, illness, and starting certain medicines can precipitate an attack
    \item \textbf{Tophi:} chronic crystal deposition forms nodules in joints, tendons, and under the skin
    \item \textbf{Kidney effects:} uric acid can also cause kidney stones and contribute to kidney disease
\end{itemize}

\section*{Clinical Features}
Gout has distinct acute attacks and chronic features.

\begin{itemize}
    \item Sudden, severe joint pain, often waking the person at night
    \item The big toe is affected in about half of first attacks
    \item Red, hot, swollen, and extremely tender joints, with pain on even light touch
    \item Attacks that peak within 12--24 hours and resolve over days to two weeks
    \item Recurrent attacks in the same or other joints
    \item Tophi: firm nodules, often at the elbows, fingers, ears, and over joints
    \item Chronic joint pain, stiffness, and deformity in long-standing gout
    \item Kidney stones causing flank pain
\end{itemize}

\section*{Differential Diagnosis}
Other causes of acute joint pain and swelling must be distinguished from gout.

\begin{itemize}
    \item Septic arthritis, a joint infection that is a medical emergency
    \item Pseudogout, caused by calcium pyrophosphate crystals rather than urate
    \item Cellulitis, a skin infection that can mimic gout in the foot
    \item Trauma and fracture
    \item Rheumatoid arthritis and psoriatic arthritis
    \item Reactive arthritis
    \item Osteoarthritis, which causes chronic rather than acute pain
    \item Bunions and other foot conditions
\end{itemize}

\section*{When to Seek Medical Care}
See a doctor for a first attack of suspected gout, to confirm the diagnosis and plan treatment, and for recurrent attacks to consider urate-lowering therapy. Anyone with a hot, swollen joint and fever should be assessed urgently, because infection must be excluded.

\textbf{Call 000 (or local emergency services) for:}
\begin{itemize}
    \item A hot, swollen, painful joint with fever, which may be septic arthritis and requires urgent treatment
    \item Sudden severe chest pain or breathlessness, which may occur with other crystal-related conditions
    \item Signs of severe allergic reaction to gout medicines
    \item Inability to use a limb or pass urine with flank pain, which may indicate a kidney stone
\end{itemize}

\section*{Diagnosis and Investigations}
Gout is diagnosed from the clinical picture, blood tests, and sometimes joint fluid analysis.

\begin{itemize}
    \item \textbf{Clinical history:} the pattern of sudden, severe attacks in typical joints supports the diagnosis
    \item \textbf{Blood uric acid:} levels are usually high, but can be normal during an acute attack
    \item \textbf{Joint aspiration:} removing fluid from the joint and examining it for urate crystals is the gold standard and also excludes infection
    \item \textbf{Imaging:} X-rays show chronic joint damage; ultrasound and DECT can detect crystal deposits
    \item \textbf{Kidney function:} tested because gout and kidney disease are linked
    \item \textbf{Response to treatment:} rapid improvement with anti-inflammatory treatment supports the diagnosis
\end{itemize}

\section*{Management}
Gout management has two parts: treating acute attacks and lowering uric acid to prevent future attacks and damage.

\begin{itemize}
    \item \textbf{Acute attacks:} non-steroidal anti-inflammatory drugs (NSAIDs), colchicine, or corticosteroids, started early and taken at full dose until the attack settles
    \item \textbf{Ice and rest:} elevating and icing the affected joint eases pain
    \item \textbf{Urate-lowering therapy:} allopurinol is the first-line medicine; febuxostat is an alternative. These are started once the acute attack has settled
    \item \textbf{Treat-to-target:} the dose is adjusted to bring blood uric acid below the target level, usually under 0.36 mmol/L
    \item \textbf{Prevention of flare on starting therapy:} low-dose colchicine or an NSAID is used for the first months of urate-lowering therapy, because starting it can trigger attacks
    \item \textbf{Lifestyle:} reducing alcohol, especially beer; limiting purine-rich foods; maintaining a healthy weight; and drinking plenty of water
    \item \textbf{Lifelong treatment:} most people with recurrent gout take urate-lowering therapy indefinitely
\end{itemize}

\section*{Complications}
\begin{itemize}
    \item Recurrent, disabling attacks
    \item Chronic gouty arthritis with joint damage and deformity
    \item Tophi causing pain, nerve compression, and skin ulceration
    \item Kidney stones and chronic kidney disease
    \item Increased risk of cardiovascular disease, which is associated with gout and its risk factors
    \item Side effects of medicines, including gastrointestinal problems with NSAIDs and colchicine
    \item Reduced mobility and quality of life
\end{itemize}

\section*{Prognosis and Outcomes}
Gout is one of the most treatable forms of arthritis. Acute attacks resolve within days to two weeks, and with effective urate-lowering therapy, attacks become less frequent and eventually stop, tophi shrink, and joint damage is prevented. Most people need lifelong medication, and adherence is important because stopping therapy leads to recurrent attacks within months. People who achieve and maintain target uric acid levels have an excellent outlook and can lead normal, pain-free lives.

\section*{Prevention and Self-Care}
\begin{itemize}
    \item Take urate-lowering medicine daily as prescribed, even when you have no symptoms
    \item Keep uric acid levels in the target range with regular blood tests
    \item Limit alcohol, particularly beer, and avoid binge drinking
    \item Limit purine-rich foods such as organ meats, red meat, and some seafood
    \item Drink plenty of water to help the kidneys excrete uric acid
    \item Maintain a healthy weight and lose weight gradually, as rapid weight loss can trigger attacks
    \item Avoid dehydration, especially around exercise and hot weather
    \item Treat acute attacks early with the medicines prescribed by your doctor
    \item Discuss all medicines with your doctor, as some, including diuretics, can raise uric acid
\end{itemize}

\section*{Sources and Further Reading}
The following reputable sources provide further information for healthcare students and patients seeking reliable, current advice about gout.

\begin{itemize}
    \item Australian Rheumatology Association. \textit{Gout information for patients.} https://rheumatology.org.au/
    \item Healthdirect Australia. \textit{Gout.} https://www.healthdirect.gov.au/
    \item NPS MedicineWise. \textit{Gout medicines: using them well.} https://www.nps.org.au/
    \item Arthritis Australia. \textit{Gout fact sheet.} https://arthritisaustralia.com.au/
    \item NHS. \textit{Gout: overview and treatment.} https://www.nhs.uk/
    \item MSD Manual. \textit{Gout.} https://www.msdmanuals.com/
\end{itemize}
